% Included within Introduction, before the present contribution.
\begingroup
\footnotesize
\setlength{\tabcolsep}{4pt}
\renewcommand{\arraystretch}{1.06}
\setlength{\LTcapwidth}{\linewidth}
\begin{longtable}{@{}>{\raggedright\arraybackslash}p{2.6cm}>{\raggedright\arraybackslash}p{4.05cm}>{\raggedright\arraybackslash}p{4.15cm}>{\raggedright\arraybackslash}p{4.85cm}@{}}
\caption{Positioning of the present study relative to constitutive learning and reduced micromechanics. Properties refer to the specified formulation and its hypotheses, not to all models sharing an architecture name.}
\label{tab:literature}\\
\toprule
Reference / formulation & Represented object & Imposed structure & Scope relevant to this work\\
\midrule
\endfirsthead
\multicolumn{4}{l}{Table \thetable\ continued}\\
\toprule
Reference / formulation & Represented object & Imposed structure & Scope relevant to this work\\
\midrule
\endhead
\bottomrule
\endfoot
As'ad et al.~\cite{asad2022}
& Scalar energy in objective Green strain; stress from differentiation
& Potential structure; ICNN convexity in Green strain
& Stability depends on stress state; strain convexity is not a global polyconvexity certificate.\\*
\addlinespace
Xu et al.~\cite{Xu2021}
& Incremental stress update using a learned lower-triangular matrix factor
& The product $\bm L\bm L^T$ is positive semidefinite, and positive definite if $\bm L$ is nonsingular
& Not necessarily an energy Hessian; the complete update derivative includes state-dependence terms.\\
\addlinespace
Klein et al.~\cite{klein2022}: invariant model
& Convex nondecreasing energy network acting on polyconvex invariants
& Objectivity and polyconvexity by construction for the selected invariants
& Feature choice restricts approximation; the lattice example illustrates why enlarging the core can be insufficient.\\
\addlinespace
Klein et al.~\cite{klein2022}: gradient model
& Network of deformation minors with material-group symmetrization
& Polyconvexity and exact selected finite-group symmetry
& Objectivity is approximate. The transverse-isotropy example also approximates the continuous symmetry by a finite subgroup.\\
\addlinespace
Linden et al.~\cite{linden2023}
& Invariant energy network with analytic growth and reference corrections
& Compatible polyconvex normalization for the stated isotropic and transverse-isotropic cases
& Global nonnegative energy is separate from zero reference energy and stress.\\
\addlinespace
Linka and Kuhl~\cite{Linka2023}
& CANN using invariant energy terms and interpretable learned coefficients
& Architecture encodes constitutive restrictions for the specified building blocks
& Model discovery is within a specified invariant family, not a completeness result for arbitrary anisotropy.\\
\addlinespace
Tac et al.~\cite{Tac2022}
& Scalar NODE flows construct monotone invariant-energy derivatives
& Integrable invariant expansion and polyconvexity under the stated nonnegative/monotone construction
& Tissue examples are incompressible with selected fiber invariants; compressible extensions are discussed.\\
\addlinespace
Ghaderi et al.~\cite{Ghaderi2020}
& Directional learning agents, microsphere summation, and maximum-history variables
& Physics-based decomposition for polymer elasticity and inelastic effects
& Includes history variables. Weight signs alone do not establish a general convexity guarantee.\\
\addlinespace
Abdolazizi et al.~\cite{Abdolazizi2025}
& CKAN scalar energy; spline edges, reference correction, and symbolic simplification
& Energy differentiation and partial monotonicity in selected inputs
& Partial monotonicity does not establish input convexity. Polyconvexity is identified as a possible extension.\\
\addlinespace
Kalina et al.~\cite{Kalina2024}
& Invariants formed with learned generalized structure tensors of orders two, four, or six
& Objective energy representation and identification of anisotropy through the tensor construction
& Learns anisotropy; polyconvexity requires additional conditions on the invariants.\\
\addlinespace
Klein et al.~\cite{klein2026}
& Objective triclinic polyconvex features, with optional finite-group symmetrization
& Polyconvexity and objectivity without requiring a prescribed nontrivial material symmetry
& Existing group-unspecified construction; the present focus is paired features with an isotropic reference derivative.\\
\addlinespace
Thakolkaran et al.~\cite{thakolkaran2025}
& Input-convex monotone spline KAN for isotropic compressible energy
& Polyconvex composition on the constrained spline domains
& Finite-secant tails require a separate join check; our integrated-hat realization uses analytic $C^2$ matching.\\
\addlinespace
Barnett et al.~\cite{barnett2023}; Ares De Parga et al.~\cite{aresdeparga2026nonlinear}
& Primary/secondary state partition; secondary coordinates learned from primary coordinates
& Low-dimensional closure integrated into a projection-based equilibrium or evolution model
& Retains a reduced solve. Direct strain substitution is a separate operating mode here.\\
\addlinespace
Farhat et al.~\cite{farhat2015}
& Fixed-weight ECSW applied to reduced nonlinear finite element dynamics
& Lagrangian structure preservation under the specified hypotheses
& Does not automatically cover state-dependent weights or independent macroscopic stress cubature.\\
\addlinespace
Hern\'andez et al.~\cite{hernandez2026}
& Nonlinear-manifold virtual work with manifold-adaptive empirical cubature
& Adaptive integration constraints and reduced equilibrium
& Supplies the adaptive cubature used here. Integration accuracy and macroscopic constitutive guarantees remain distinct.\\
\end{longtable}
\endgroup
